\subsection{Scope and interpretive boundaries}
Before presenting the experimental results, three scope limitations
should be noted to avoid over-interpretation.  First, the learned
edge relevance and surface relevance $C_j$ reflect
response-consistent interaction patterns under the forecasting task;
they are \emph{not} direct measurements of physical contact force,
contact pressure, or jamming probability.  Second, the case study
uses a single tunnel section with 49 excavation steps (44 effective
sequence samples, 12 in the stratified test set); all quantitative
indicators should be read as illustrative rather than statistically
conclusive.  Third, the spatial consistency indicators (Moran's $I$,
component CV, attention--geology correlation) quantify the internal
and external plausibility of the learned representation, not its
engineering correctness.  These boundaries are revisited in the
Discussion.

\subsection{Experimental setup}
\label{sec:experimental_setup}
All structural and training hyperparameters follow the Implementation
details in Section~\ref{sec:methods} (summarised here for readability):
historical window $K=5$, forecasting horizon $h=1$; candidate edge
distance threshold $\tau_{\mathrm{edge}}=2.0$\,m; signed-normal
compatibility threshold $\eta_{\min}=0.3$; active-zone radial extent
$\tau_{\mathrm{zone}}=5.0$\,m; TBM surface resolution 0.5\,m (1352 nodes);
GNN encoder with 2 layers, hidden dimension 32, dropout 0.3; GRU hidden
dimension 64, 1 layer; AdamW optimiser, learning rate $10^{-3}$, weight
decay $10^{-2}$, batch size 32, early stopping patience 15.  To control
for geological distribution shift between train and test, samples are
partitioned using stratified random sampling based on TSP shear wave
velocity ($V_s$) quartiles, yielding four geological strata each split
70/15/15 into train/val/test.  All models are evaluated on the
held-out stratified test set using MAE, RMSE, $R^{2}$, and Pearson
correlation.

\subsection{Prediction performance}
Table~\ref{tab:prediction} reports the global prediction performance
(Figure~\ref{fig:prediction}).  The proposed framework achieves MAE of
0.417 and Pearson $r$ of 0.567, outperforming all baselines on MAE and
achieving the only positive $R^{2}$ (0.133).  The stratified split
controls for geological heterogeneity between train and test, allowing
the geometric prior and learned rock--machine interactions to
generalise; under the previous mileage split, all models exhibited
negative $R^{2}$ due to a severe distribution shift in $V_s$ between
the training and test intervals (shift $= 2.31$ standard deviations,
Kolmogorov--Smirnov $p < 0.001$).  The proposed framework's primary
contribution lies in providing spatially explicit
interaction structure that monitoring-only models cannot access---a
trade-off examined in Section~\ref{sec:tradeoff}.

\begin{table}[htbp]
\centering
\caption{Global prediction performance on the stratified test set
($n=12$ samples drawn from four geological strata defined by TSP shear
wave velocity quartiles, ensuring the test distribution matches the
training distribution).  The proposed framework achieves the lowest MAE
and the only positive $R^{2}$ among all models, indicating that the
geometric prior and learned rock--machine interactions generalise when
geological heterogeneity between train and test is controlled.}
\label{tab:prediction}
\begin{tabular}{lcccc}
\toprule
Method & MAE & RMSE & $R^{2}$ & Pearson $r$ \\
\midrule
Persistence & 0.465 & 0.590 & $-0.008$ & 0.600 \\
XGBoost + TSP & 0.455 & 0.606 & $-0.114$ & 0.496 \\
LSTM & 0.488 & 0.602 & $-0.109$ & 0.457 \\
TSP-LSTM & 0.487 & 0.619 & $-0.174$ & 0.449 \\
Proposed framework & \textbf{0.417} & \textbf{0.540} & $\mathbf{0.133}$ & 0.567 \\
\bottomrule
\end{tabular}
\end{table}

\begin{figure}[htbp]
\centering
\includegraphics[width=0.9\linewidth]{fig3_prediction_comparison.pdf}
\caption{Predicted versus observed standardised monitoring responses
for five target variables on the stratified test set.  Each panel
corresponds to one response variable: (a)~advance rate, (b)~torque,
(c)~thrust, (d)~penetration, (e)~shield pressure.  Different curves
represent the proposed framework and baseline models; observed values
are shown as solid lines, predictions as dashed/dotted lines.  MAE
values in the legend are variable-specific values for the displayed
panel and should not be confused with the global multi-variable metrics
in Table~\ref{tab:prediction}.  Under the stratified split, the proposed
framework achieves positive $R^{2}$ and outperforms all baselines on
MAE, confirming that the geometric prior contributes predictive signal
when geological heterogeneity between train and test is controlled.}
\label{fig:prediction}
\end{figure}

\subsection{Structural ablation}
Table~\ref{tab:ablation} and Figure~\ref{fig:ablation} present the structural ablation results.  Four
ablation variants isolate the contribution of each design component.

\begin{table}[htbp]
\centering
\caption{Structural ablation results on the stratified test set.  Each
row removes or replaces one design component from the full framework.
The geometry-only baseline is excluded from this table because it has no
trainable parameters and does not produce predictions; its spatial
consistency is reported in Table~\ref{tab:spatial_consistency}.  Under
the stratified split, the full framework achieves the lowest MAE,
confirming that the geometric prior contributes predictive signal
rather than merely imposing regularisation.}
\label{tab:ablation}
\begin{tabular}{lcccc}
\toprule
Variant & MAE & RMSE & $R^{2}$ & Pearson $r$ \\
\midrule
Full framework & \textbf{0.417} & \textbf{0.540} & $\mathbf{0.133}$ & 0.567 \\
Without monitoring input & 0.558 & 0.666 & $-0.540$ & $-0.141$ \\
Randomised rock--machine edges & 0.462 & 0.598 & $-0.205$ & 0.438 \\
Without geometric prior ($\beta=0$) & 0.539 & 0.635 & $-0.319$ & $-0.012$ \\
\bottomrule
\end{tabular}
\end{table}

\begin{figure}[htbp]
\centering
\includegraphics[width=0.9\linewidth]{fig4_ablation_study.pdf}
\caption{Structural ablation comparison.  (a)~Prediction metrics (MAE
and $R^{2}$) across four ablation variants; (b)~illustrative spatial
diagnostics generated during figure preparation.  The traceable
checkpoint-based spatial evidence used for interpretation is reported in
Table~\ref{tab:spatial_consistency}.  The geometry-only baseline is
omitted from panel~(a) because it has no trainable parameters and does
not produce predictions.  Under the stratified split, the full framework
achieves the lowest MAE and the only positive $R^{2}$, indicating that
the geometric prior contributes predictive signal rather than merely
imposing regularisation.}
\label{fig:ablation}
\end{figure}

The \emph{No Monitoring} ablation removes the monitoring vector from the GRU
input.  Pearson $r$ collapses from 0.567 to $-0.141$, suggesting that
temporal monitoring dynamics carry essential predictive signal that the
static graph structure alone cannot replace.

The \emph{Randomised Edges} ablation replaces geometry-constrained
rock--machine edges with random edges of the same cardinality.  Pearson $r$
drops from 0.567 to 0.438, suggesting that the spatial structure of
candidate relations---not merely the presence of cross-type
edges---contributes to learning meaningful interaction patterns.

The \emph{No Geometric Prior} ablation sets $\beta=0$.  Under the
stratified split, this variant achieves higher MAE (0.539) and near-zero
Pearson $r$ ($-0.012$) compared with the full framework, indicating that the
geometric prior contributes predictive signal rather than merely
imposing regularisation.  This contrasts with the mileage split, where
removing the prior appeared to improve prediction; the stratified split
isolates the prior's contribution by controlling for geological
heterogeneity between train and test.  Spatial consistency is examined
separately below using post-hoc evidence regenerated from the saved
checkpoint.

The \emph{Geometry-only baseline} replaces all learned edge scores with
the fixed geometric prior $g_{ij}=\exp(-d_{ij}/\tau_{\mathrm{edge}})$
(no trainable parameters).  Because it has no learned forecasting head,
it does not produce predictions and is therefore excluded from
Table~\ref{tab:ablation}; it is, however, retained in the spatial
consistency comparison (Table~\ref{tab:spatial_consistency}) because it
yields a well-defined surface relevance map $C_j$ directly from the
geometric prior.  This baseline serves as a non-learned reference point:
it provides a fixed non-learned spatial reference for evaluating whether
the learned relevance refines or simply mirrors the candidate-edge geometry.

\subsection{Spatial consistency of learned attention}
\label{sec:spatial_consistency}
A central claim of this framework is that geometry-constrained graph
construction produces spatially consistent attention patterns that align
with engineering expectations.  We evaluate this claim using three
quantitative spatial consistency indicators, each targeting a different
aspect of spatial validity.

\paragraph{Attention--geology correlation.}
We quantify whether the mean surface relevance
$\bar{C}(c_k)=|V^m|^{-1}\sum_j C_j(c_k)$ co-varies with the TSP P-wave
velocity interpolated at the same test chainages.  This check is treated
as an external plausibility indicator, not as a required success
criterion: in a response-supervised setting, learned relevance may also
reflect machine operating state and temporal response history.  The
post-hoc checkpoint evidence shows only weak attention--geology alignment
for the learned model (Pearson $r=0.067$, $p=0.837$; Spearman
$r=0.294$, $p=0.354$), whereas the geometry-only prior has a stronger
positive velocity correlation (Pearson $r=0.623$, $p=0.031$).  Thus the
current checkpoint should not be interpreted as proving that low-velocity
geological zones directly receive higher learned relevance.

\paragraph{Spatial autocorrelation of surface relevance.}
Spatially consistent relevance should exhibit positive spatial
autocorrelation: nearby TBM surface nodes should have similar $C_j$ values,
rather than random spatial noise.  We compute Moran's $I$
\citep{moran1950notes} for the surface relevance values on the TBM surface
at each test chainage:
\[
I(c) = \frac{n}{S_0} \cdot
\frac{\sum_i \sum_j w_{ij}\,(C_i(c) - \bar{C}(c))\,(C_j(c) - \bar{C}(c))}
     {\sum_i (C_i(c) - \bar{C}(c))^2},
\]
where $n = |V^m|$ is the number of TBM surface nodes, $w_{ij}$ is the
spatial weight between nodes $i$ and $j$ (binary $k$-nearest-neighbour,
$k=8$, row-standardised), and $S_0 = \sum_i \sum_j w_{ij}$.  Moran's $I$
ranges from approximately $-1$ (perfect dispersion) through $0$ (spatial
randomness) to $+1$ (perfect spatial clustering).  The reported value is
the mean $I(c)$ across all test chainages.  The full framework achieves
Moran's $I=0.925$, indicating strongly clustered surface relevance.  The
geometry-only baseline yields a lower and less stable mean Moran's $I$
(0.170), suggesting that the learned scores are not merely a direct copy
of the fixed geometric prior.

\paragraph{Component-level relevance differentiation.}
If the surface relevance captures meaningful rock--machine interaction
structure, different TBM components (cutterhead, front/middle/tail shield)
should exhibit distinct mean relevance levels, reflecting their different
interaction intensities with the surrounding rock.  We measure the
coefficient of variation (CV) of component-mean relevance across the four
components:
\[
\mathrm{CV} = \frac{\sigma_{\mathrm{comp}}}{\mu_{\mathrm{comp}}},
\quad
\mu_{\mathrm{comp}} = \frac{1}{4}\sum_{k=1}^{4} \bar{C}_k,
\quad
\bar{C}_k = \frac{1}{|V_k^m|}\sum_{j \in V_k^m} C_j,
\]
where $V_k^m$ denotes the set of TBM surface nodes belonging to component
$k$, and $\sigma_{\mathrm{comp}}$ is the standard deviation of the four
component means.  CV is dimensionless and scale-invariant, making it
suitable for comparing differentiation across models with different
absolute relevance magnitudes.  The full framework achieves mean
CV\,$=0.193$, higher than the geometry-only baseline (0.105), indicating
moderate component-level differentiation.

\paragraph{Degree-control check.}
Because rock--machine candidate edges are spatially constrained, nodes
with more incident edges could appear more relevant simply because of
edge count.  We therefore correlate degree-normalised relevance $C_j$
with incident candidate-edge degree $d_j$ across TBM surface nodes.  The
mean Pearson correlation is small ($r=0.043$), suggesting that the
degree-normalised definition of $C_j$ largely removes edge-count inflation.

\begin{table}[htbp]
\centering
\caption{Post-hoc spatial evidence regenerated from the saved checkpoint.
Higher Moran's $I$ and component CV indicate greater spatial structure.
Rel.--Geo.\ $r$ is the Pearson correlation between mean surface relevance
and TSP P-wave velocity across test chainages.  Degree-control $r$
correlates $C_j$ with incident candidate-edge degree and should be close
to zero if relevance is not merely edge-count driven.}
\label{tab:spatial_consistency}
\begin{tabular}{lcccc}
\toprule
Variant & Rel.--Geo.\ $r$ & Moran's $I$ & CV & Degree $r$ \\
\midrule
Full framework & 0.067 & 0.925 & 0.193 & 0.043 \\
Geometry-only baseline & 0.623 & 0.170 & 0.105 & -- \\
\bottomrule
\end{tabular}
\end{table}

These indicators support a narrower but better-evidenced interpretation:
the saved model produces spatially clustered and component-differentiated
surface relevance, and this relevance is not simply explained by the
number of incident candidate edges.  The current evidence does not support
a strong claim that learned relevance is monotonically higher in
low-velocity geological zones; this remains a target for additional
stability and sensitivity experiments.

\subsection{Spatial interpretation outputs}
Figure~\ref{fig:hotspot} presents the response-consistent hotspot maps
projected onto the TBM shield surface at selected chainage positions.
Colour intensity indicates the aggregated edge relevance incident to each
surface node, revealing which shield regions the model identifies as most
strongly coupled to the surrounding rock mass.

\begin{figure}[htbp]
\centering
\includegraphics[width=0.95\linewidth]{fig5_hotspot_maps.pdf}
\caption{Response-consistent surface relevance hotspots on the
unwrapped TBM surface at three representative test chainages.
The colour indicates degree-normalised surface relevance $C_j$,
rescaled within each chainage for visualisation only; different
chainages are therefore not directly comparable in absolute colour
intensity.  Vertical dashed lines separate cutterhead, front shield,
middle shield, and tail shield regions; horizontal dotted lines mark
crown, springline, and invert positions.  (a)~Chainage 41\,m (early
test set); (b)~chainage 45\,m (mid test set); (c)~chainage 48\,m
(late test set).  The maps show where the learned rock--machine
interaction relevance concentrates on the TBM surface.}
\label{fig:hotspot}
\end{figure}

Figure~\ref{fig:chainage_evolution} shows the chainage evolution of
component-wise surface relevance alongside the TSP velocity field and monitoring
responses.  Panel~(a) displays $C_j$ decomposed by TBM component
(cutterhead, front shield, middle shield, tail shield) at each test
chainage, with the cross-component mean overlaid as a dashed line;
panel~(b) shows the TSP P-wave velocity profile with low-velocity
anomalies (below the 25th percentile) highlighted in red; panels~(c)--(g)
show the five monitoring response variables.

\begin{figure}[htbp]
\centering
\includegraphics[width=0.9\linewidth]{fig6_chainage_evolution.pdf}
\caption{Chainage evolution of component-wise surface relevance,
geological context, and monitoring response.  (a)~Surface relevance
$C_j$ decomposed by TBM component (cutterhead, front shield, middle
shield, tail shield; degree-normalised, unstandardised) at each test
chainage, with the cross-component mean overlaid as a dashed line;
(b)~TSP P-wave velocity profile with low-velocity anomalies (below
25th percentile) highlighted in red; (c--g)~monitoring response
variables (advance rate, torque, thrust, penetration, shield pressure).
The co-evolution of component-wise surface relevance with geological
context and monitoring response reveals how interaction emphasis
migrates across TBM components as the TBM advances through different
geological zones.}
\label{fig:chainage_evolution}
\end{figure}

The hotspot maps reveal spatially heterogeneous interaction patterns: the
front shield and cutterhead-proximal region consistently exhibit higher
relevance than the tail shield, consistent with the engineering expectation
that rock--machine interaction is concentrated near the excavation face.
The chainage evolution view shows that interaction relevance shifts across
components as the TBM advances, suggesting that the model captures
geologically meaningful spatial variation in rock--machine coupling.  These
spatial patterns are not recoverable from monitoring-only sequence models,
which lack any spatial representation of the interaction geometry.

Figure~\ref{fig:decision_support} illustrates how the spatial
interpretation outputs \emph{could} inform excavation decision-making by
integrating predicted responses, attention-based signals, and geological
context into a unified spatial view.  This is a conceptual demonstration,
not a validated decision-support tool.

\begin{figure}[htbp]
\centering
\includegraphics[width=0.9\linewidth]{fig7_decision_support.pdf}
\caption{Illustrative decision-support scenario on the test set
(chainage 41--48\,m).  (a)~Predicted vs.\ observed advance rate with
geological scenario annotation (low-velocity zone, transition zone,
normal zone); (b)~surface relevance signal (mean $C_j$) with TSP
velocity overlay on the right axis; the red dashed line marks the
alert threshold (75th percentile of $C_j$), and the shaded red region
highlights chainages where $C_j$ exceeds this threshold;
(c)~composite attention--error indicator combining min--max normalised
surface relevance and normalised prediction error.  This indicator is a
heuristic cue for illustration only; it is not a calibrated risk metric
and has not been validated for operational use.  The spatial reasoning
chain---linking geological context, attention-based interaction
assessment, and prediction error---is \emph{facilitated by} (not
exclusive to) the graph-based spatial representation.  Further
validation on larger datasets is needed before any operational
deployment.}
\label{fig:decision_support}
\end{figure}

\subsection{Robustness and sanity checks}
\label{sec:robustness}
Table~\ref{tab:robustness} summarises the checks that are currently
backed by regenerated evidence files.  The reported values are produced
by reloading the saved checkpoint, rebuilding the same stratified test
set, and recomputing prediction and spatial evidence.  Additional
multi-seed, label-shuffle, and threshold-sensitivity experiments are
defined in the reproducibility protocol but are not reported here until
their output files are generated.

\begin{table}[htbp]
\centering
\caption{Traceable robustness and sanity checks regenerated from the
saved checkpoint.  The evidence source is
\texttt{outputs/evidence/posthoc\_evidence.json}.}
\label{tab:robustness}
\begin{tabular}{llc}
\toprule
Check & Metric & Result \\
\midrule
Checkpoint replay & MAE & 0.417 \\
Checkpoint replay & RMSE & 0.540 \\
Checkpoint replay & $R^2$ & 0.133 \\
Surface clustering & mean Moran's $I$ & 0.925 \\
Component differentiation & mean CV & 0.193 \\
Edge-count control & mean Pearson $r(C_j,d_j)$ & 0.043 \\
Geometry-only reference & mean Moran's $I$ & 0.170 \\
\bottomrule
\end{tabular}
\end{table}

\emph{Checkpoint replay} verifies that the saved graph-sequence model
reproduces the reported prediction metrics when the test set is rebuilt
from the configuration file.  The \emph{edge-count control} rules out the
possibility that higher $C_j$ simply reflects a higher number of incident
candidate edges.  The degree-normalised definition
$C_j = \frac{1}{d_j+\epsilon}\sum_i \mathrm{softplus}(s_{ij}^{rm})$
largely decouples relevance from edge count in the current checkpoint.
Seed stability, label permutation, and candidate-edge threshold
sensitivity remain required follow-up checks before making stronger
claims about robustness across training initialisations or graph
construction parameters.
